\chapter{Requirements Analysis}
\label{ch:requirements}

\section{Stakeholders}

The system is designed for three kinds of stakeholder. \textbf{Researchers}
(the project team, and future students building on this corpus) need an
annotated dataset and a classifier whose reported numbers can be trusted --
i.e., evaluated against an honest ceiling rather than an assumed perfect
one. \textbf{Nonprofit fundraising practitioners} are the eventual downstream
users of the strategy-to-outcome findings: which strategies are worth
training callers to use. \textbf{Dialogue-system builders} are the users of
the donation-intent classifier and its generic-domain extension, who need a
model that transfers to a related but distinct persuasive setting rather than
one that has memorized this specific corpus.

\section{Functional Requirements}

\begin{description}
\item[FR1 -- Multi-label strategy taxonomy.] The system shall define a
  fixed, mutually-understood taxonomy of persuasion strategies, organized
  into higher-level categories, sufficient to describe the range of tactics
  observed in the corpus.
\item[FR2 -- Exhaustive turn-level annotation.] The system shall label every
  persuader turn in the annotated portion of the corpus with the
  \emph{complete set} of strategies it uses (possibly the empty set), not a
  single dominant strategy, using the preceding dialogue turns as context
  where the turn's meaning depends on it.
\item[FR3 -- Automatic strategy prediction.] The system shall, given a
  persuader turn (and up to five preceding turns of dialogue context),
  predict the probability of each strategy/category being present, and
  convert these probabilities into a discrete label set via a
  decision procedure fitted on held-out data.
\item[FR4 -- Dual objective support.] The system shall support being
  evaluated and operated at two distinct, independently useful operating
  points: one that maximizes micro-F1 (overall labeling accuracy, dominated
  by common strategies) and one that maximizes macro-F1 (equal weight per
  strategy, sensitive to rare-strategy performance), since a fundraising
  application interested in \emph{rare but high-value} tactics needs the
  latter, and one interested in overall labeling throughput needs the
  former.
\item[FR5 -- Donation-outcome estimation.] The system shall, given the
  (gold) strategy composition of a dialogue, estimate the strategy set's
  association with the dialogue's binary donation outcome, report effect
  sizes (odds ratios) with significance levels corrected for multiple
  comparisons, and report the model's out-of-sample explanatory power.
\item[FR6 -- Donation-intent classification.] The system shall, given a full
  persuader--persuadee dialogue, classify whether it ends in a donation
  intent and, if so, of what type (outright, deferred, or conditional).
\item[FR7 -- Cross-domain transfer check.] The system's dialogue-classification
  architecture shall be evaluated on a second, structurally similar but
  topically distinct persuasive-dialogue domain to test whether its design
  generalizes or has overfit to charitable solicitation specifically.
\end{description}

\section{Non-Functional Requirements}

\begin{description}
\item[NFR1 -- Evaluated against an honest ceiling.] Every classification
  result shall be reported alongside a measured annotation-reliability
  ceiling (Section~\ref{sec:ceiling}) -- how often independent labelings of
  near-duplicate turns agree with each other -- rather than compared only to
  a theoretical maximum of 1.0, which no annotation process, human or
  automated, can be expected to reach.
\item[NFR2 -- No leakage across the train/validation/test split.] Because
  turns from the same dialogue share context and style, the split into
  train/validation/test shall be performed at the \emph{dialogue} level, not
  the turn level, and every downstream statistic (cross-fitted decision
  thresholds, bootstrap confidence intervals, significance tests) shall
  resample or fold by dialogue for the same reason.
\item[NFR3 -- Rater-agnostic modeling.] The classifier shall not be given
  the identity of the human/LLM annotator who produced a turn's label as an
  input feature, even though doing so measurably increases apparent accuracy,
  because a production system will not know who \emph{would have} labeled an
  unseen turn, and training on annotator identity would be learning a
  shortcut rather than the underlying task.
\item[NFR4 -- Reproducibility.] All randomized procedures (data splitting,
  weight initialization, data augmentation, cross-fitted decision thresholds,
  bootstrap resampling) shall use fixed seeds, and a subset of ensemble
  members (\emph{sentinels}) shall be kept byte-for-byte identical across
  consecutive design iterations specifically so that any change in their
  reported score signals a bug rather than a real effect
  (Section~\ref{sec:design-methodology}).
\item[NFR5 -- Compute budget.] The full training-and-evaluation pipeline
  shall complete within a single free-tier Kaggle GPU session (an NVIDIA
  T4$\times$2 instance, walltime limit on the order of hours, weekly quota
  limits), and shall be validated on a small data subsample (a
  \emph{smoke test}) before any full run is committed, to avoid spending
  scarce GPU-hours on a run that fails partway through from a code defect.
\item[NFR6 -- Statistical honesty.] Confidence intervals shall be computed
  by cluster bootstrap over dialogues (not turns), and any comparison
  between two models or two design choices shall be accompanied by a paired
  significance test rather than a bare point-estimate difference.
\end{description}

\section{Data Requirements}

The corpus provides 10{,}600 persuader turns across 1{,}017 dialogues; the
taxonomy requires 41 strategy labels organized into 11 categories
(Section~\ref{sec:taxonomy}); the annotation process is required to report
its own inter-annotator (and, where an LLM is involved, inter-run) agreement
so that classifier performance can be judged against it rather than in a
vacuum (Section~\ref{sec:ceiling}).

\section{Tooling and Environment Requirements}

Model development targets PyTorch with the HuggingFace \texttt{transformers}
library for encoder fine-tuning, \texttt{scikit-learn} for shallow baselines,
thresholding, and calibration, and \texttt{statsmodels} for the logistic
outcome models in Phase~4. Training runs on Kaggle's free GPU tier; because a
single account's weekly GPU quota is not guaranteed to survive a full
seventeen-iteration research process, redundant training across a second
Kaggle account was adopted as an infrastructure requirement, not a
convenience (Section~\ref{sec:design-methodology}).
